\section{Mechanisms}

The four mechanisms answer four different questions: where to look, what the
failure means, when to stop, and what the run cost. Each is stated with the
observation that would falsify it, because the parts that were easiest to
describe were the parts most likely to be doing nothing.

\subsection{The locator card}
\label{sec:locator}

\paragraph{Motivation.} The long-horizon runs did not fail for lack of
information; they failed while reading. In the baseline the agent reopened the
same file repeatedly, never wrote a patch, and exhausted its rounds. It lacked
position: the failing test names a module, the seeded frame names a function
that merely calls the culprit, and the function to edit sits a short distance
away in the module's reference graph.

\paragraph{Approach.} Preflight yields the frames of the failing test with no
model call. One-hop entries are the symbols those frames name; two-hop entries
come from a breadth-first walk over a graph built from the module's abstract
syntax tree, with top-level functions, classes and module-level assignments as
nodes and name references inside bodies as edges. Candidates are ranked by task
words in the symbol name, then match strength, hop count and definition line.
The first three become cards carrying path, symbol, line span and a reason from
a closed set.

\begin{verbatim}
entries  <- symbols named by the failing frames
frontier <- entries
repeat at most twice:
    add top-level symbols referenced from frontier (in-module AST)
rank by (-task words in name, match tier, hop, line)
emit the first three entries as cards
\end{verbatim}

\paragraph{Criterion.} The card is judged by the first read. When it was first
delivered, the first read landed in the golden region on 4 of 4\fact{E6}
instances, and 1 of 4\fact{E6} resolved --- a mechanism check on four instances
of one task. The card names locations and never assertions: no test name, no
expected value, so it can point at the right function without handing over the
answer.

\paragraph{Failure modes.} The walk stays inside one module: cross-file calls
are not chased, types are not resolved, dynamic dispatch is invisible, and a
culprit reachable only through a package alias degrades the ranking to
structural order. The sharpest failure is delivery: for a sequence of arms the
card was computed and scheduled, yet a search of the stored request bodies for
the symbol it named found nothing\fact{E9}. A correct card that never
arrives is indistinguishable, from the outside, from a wrong one.

\subsection{Failure-test evidence}

\paragraph{Motivation.} An issue report describes a symptom; a test states a
contract. For the tabulate formatting issue the contract is a literal expected
string in which whitespace and quoting are part of the specification, so a patch
that satisfies the prose can still fail the assertion. Passing the whole test
file is the wrong repair: the window fills with unrelated cases, and the
editable file invites fixing the test instead of the code.

\paragraph{Approach.} The raw runner output is parsed for the innermost
test-side frame, and that line is mapped with the parser to the enclosing test
function, so the unit is the full definition rather than a fixed window. The
failure excerpt from the error block is appended and bounded. The result enters
as one ordinary tool-result section under the existing share: no extra budget,
no new kind. The ledger keeps the path, the
symbol, the span and an artifact reference; the text stays in the object store.
If preflight is green nothing is emitted, and the section reaches the model only
if the assembler placed it, so a section evicted by the share disappears from
the request rather than only from the ledger.

\paragraph{Criterion.} Containment. On the real instance the unit held the
failing function and its assertion excerpt and none of the other tests in the
same file; a unit test asserts their absence from both fields, and a switch
disables the injection for a separate measurement.

\paragraph{Failure modes.} Import and collection errors carry no assertion, so
only the frame survives. The culprit may be a helper one call below the test:
the locator reaches it, the evidence section does not, and the two mechanisms
disagree. A patch can be fitted to the literal and still regress tests
elsewhere, which is why acceptance is node-level. And the excerpt is raw runner
text, so a format change yields an empty payload that looks exactly like a task
with nothing to inject.

\subsection{Mechanical gates}

\paragraph{Motivation.} The gates answer a different question from context: not
what the model knows, but when the episode ends. Given the card and the failing
test, the agent often had what it needed and still did not stop. Instructions we
could write down were about behavior, and behavior is what a prompt is worst at
changing.

\paragraph{Prompt versus mechanism.} We first wrote the required process as six
constraints --- deliver a patch early, do not re-read a span, note a conclusion
before acting, stay in scope, hand in rather than test, keep output short ---
and assembled them as resident project-rule sections, ordered after the charter
and before the task, rendered into the system message. This is the obvious
implementation, and it did nothing: 0/3\fact{E5}. Per-turn output grew rather
than shrank, as though the rules had been read as a description of the
situation rather than a procedure.

Moving the same requirement into the loop changed behavior but not outcome.
From the seventh round, if no patch had been submitted, the runtime appended a
user-message reminder on every turn; patches appeared, but the model revised
instead of submitting: 0/2\fact{E7}.

Only when the runtime took the decision away did the task resolve. The step
allows at most two patch versions; on the second it closes, the last version
goes to the verifier, and the model is not asked whether it is done. Same task,
same model: 5/5\fact{E8}.

\paragraph{The symptom trap.} The pattern is not specific to rules. Removing
visible symptoms one at a time did not move the outcome: deduplicating repeated
reads 0/3\fact{E2}, enlarging the share available to file reads 0/3\fact{E3},
raising the token ceiling 0/2\fact{E4}. Each cut tidied the transcript and left
an episode that could not terminate; the starting point, with neither card nor
gate, was 0/3\fact{E1}.

\paragraph{Criterion.} The gate is judged by termination, not compliance: did
the episode close with a decision the verifier could judge. Compliance is
unobservable in a transcript the model writes; termination is an event in the
ledger, with a timestamp and a patch attached.

\paragraph{Failure modes.} The two-version cap can end an episode whose first
direction was wrong; in this sample it did not, but it is a designed way to
fail. The gate does nothing for a model that already converges: the DeepSeek
runs submit early enough that the reminder never fires.

\subsection{Cost accounting}

\paragraph{Motivation.} The ledger used to record one number per call: total
tokens. Providers bill three things differently --- fresh input, cached input,
and output, mostly invisible thinking --- and one total makes the cheap part
and the expensive part indistinguishable. Two agents with different caching
then look comparable when they are not.

\paragraph{Approach.} We record the provider's usage fields verbatim: prompt,
completion, cached prompt, reasoning. Absent fields are recorded as zero, never
inferred; the old total keeps its meaning, so earlier ledgers replay. Because
the rendered request body is stored per call, we can also replay the bodies with
a minimal output cap, read the real prompt size and cache hit of every turn, and
subtract to recover completion.

\paragraph{What the attribution shows.} The gap is in output, not in input:
average output is 73,051\fact{D2} tokens against 12,495\fact{D2} for mcode,
while fresh input is 21,941\fact{D3} against 52,996\fact{D3}. On the same six
tasks with the same model and endpoint the real bill is \$0.0486\fact{D1} per
task against \$0.0292\fact{D1}, a ratio of 1.66\fact{D1}. The familiar claim
that we are 3\fact{H1} times more expensive is the uncached accounting, which
charges every input token as fresh; failures also concentrate the spend, with
successful tasks averaging 78,844\fact{D6} tokens and two failed tasks
accounting for 70.9\%\fact{D6} of all output.

\paragraph{Criterion.} A cost claim must be recomputable from the ledger and a
public price list; when two conventions disagree, both are printed next to the
claim. The lever that follows is stopping, not compressing.

% FACTS-USED: E1, E2, E3, E4, E5, E6, E7, E8, E9, D1, D2, D3, D6, H1
